\documentclass[12pt]{article}

% House style preamble
\input{.house-style/preamble.tex}

% Update PDF metadata
\hypersetup{
  pdftitle={Function-first without function-in-syntax: projectibility, field-relative kinds, and the limits of Thomasson's deflation},
  pdfkeywords={conceptual engineering, natural kinds, projectibility, deflationism, easy ontology, modal normativism}
}

% Working title -- provisional, confirm with Brett.
\title{Function-first without function-in-syntax\\
{\large Projectibility, field-relative kinds, and the limits of Thomasson's deflation}}
\author{Brett Reynolds \orcidlink{0000-0003-0073-7195}%
\thanks{Contact: \href{mailto:brett.reynolds@humber.ca}{brett.reynolds@humber.ca}}\\
Humber Polytechnic \& University of Toronto}
\date{\today}

\begin{document}
\maketitle

\begin{abstract}
Amie Thomasson's metaphysics is function-first: existence and modal questions turn on semantic rules for terms, while category questions turn on the roles that concepts and categories play in inquiry. I argue that this program doesn't depend on systemic-functional linguistics, and that ordinary linguistic practice is one field among others, not a neutral court over metaphysics. Thomasson leaves a realist opponent unmet: the anti-essentialist tradition built around projectibility, plural categorization, and induction. Projectibility-first realism accepts that pluralism while adding support grades, failure conditions, and demotion rules. Once a field and projection are specified, category profiles can be better or worse supported by the world. That support is neither a field-neutral essence nor reducible to conceptual usefulness.
\end{abstract}

\noindent\textbf{Keywords:} conceptual engineering; natural kinds; projectibility; deflationism; easy ontology; modal normativism

\section{Introduction}

Thomasson's work gives metaphysics a functional turn. Existence questions are handled through application conditions for terms \citep{thomasson2014ontology}. Metaphysical modal claims are object-language expressions of semantic rules rather than descriptions of modal features of reality \citep{thomasson2020norms}. Conceptual engineering becomes a central task for metaphysics once we ask which concepts and categories should be retained, revised, or introduced \citep{thomasson2025rethink}. These claims don't eliminate metaphysics. They redirect it.

Here \term{function} means role in a practice, not causal mechanism, biological purpose, or syntactic slot. At the semantic level, a term's function is fixed by rules for correct application and inference. At the field-inferential level, a category's function is fixed by the projections it licenses, the questions it helps answer, and the errors or demotions that would defeat it. The paper uses \term{function} at the level of use and inquiry, not as a claim about linguistic architecture.

A field-relative view also changes how ordinary language enters the argument. Ordinary linguistic practice is itself a field, with its own uptake patterns, competence norms, repair routines, and defeaters. Metaphysics is another field, including hard ontology as actually practiced. Thomasson often lets ordinary use discipline metaphysical practice. That may be right when metaphysicians claim to use ordinary terms, but it isn't automatic.

This paper accepts Thomasson's redirection. Its first claim is that her function-first metaphysics needs less linguistic machinery than she sometimes suggests. The philosophical result can be extracted from systemic-functional linguistics. It depends on a more general point: different parts of discourse have different semantic, inferential, and sorting roles, and metaphysics often goes wrong by treating all of them as attempts to describe deep structure.

The second claim concerns natural kinds. Thomasson is at her strongest against Kripke-Putnam essentialism about natural kinds. She has the resources to resist the thought that metaphysical inquiry discovers a privileged structure that any adequate category scheme has to mirror. But that isn't the only realist view of kinds. A large anti-essentialist tradition treats kinds through projectibility, causal support, plural categorization, and defeasible inductive use. This tradition is the unmet opponent.

The constructive alternative is a field-relative account of kind claims. A category has standing within a field when membership in it supports reliable projections: observing some features licenses expectations about others in ways that survive appropriate tests \citep{Goodman1955}. A kind claim then says that this standing is strong enough for kindhood in that field. This account grants Thomasson's anti-essentialism and her functional pluralism. It denies that those concessions settle the realist question. The disagreement isn't realism against pluralism. It's conceptual-engineering pluralism against projectibility-constrained pluralism.

\section{Thomasson's function-first metaphysics}

In \emph{Ontology Made Easy}, Thomasson explains existence questions through application conditions. These are semantic rules of use rather than definitions speakers have to state. Nor are they merely evidence conditions. They specify when applying a term would be correct, not simply when we'd be warranted in thinking the term applies \citep[93--95]{thomasson2014ontology}.

This machinery already accommodates more empirical deference than a simple conventionalist reading would suggest. Thomasson's \mention{platypus} example explicitly allows natural-kind terms to have world-deferential application conditions. Such terms can defer to the world, and to experts with specialized knowledge of the world, to determine what distinguishes one biological kind from another \citep[94]{thomasson2014ontology}. So the objection can't be that Thomasson has no room for natural-kind terms, causal discovery, or expert deference.

\emph{Norms and Necessity} extends the same functional strategy to modal discourse. At the semantic level, metaphysical modal claims express rules of use rather than tracking independent modal features. Thomasson also separates that claim from causal, nomological, dispositional, and empirical counterfactual claims. At the field-inferential level, scientific laws can guide empirical inference, and natural-kind categories can be assessed by how well they support predictive and explanatory generalizations \citep[121--123]{thomasson2020norms}.

\emph{Rethinking Metaphysics} makes the positive program explicit. Conceptual engineering isn't arbitrary when it responds to empirical, technological, and social constraints. Thomasson's discussion of \mention{fish} is especially important. Biological concepts can change as empirical discoveries and theoretical purposes change; the resulting categories can be assessed by how well they support generalization and prediction \citep[216--217]{thomasson2025rethink}.

These passages create the central convergence. Thomasson already has a function-relative, empirically constrained account of concept and category choice. The question is whether that account exhausts what a realist projectibility theory wants from kinds.

Projectibility isn't just another purpose a concept might serve. It's a way of testing whether a category's success outruns our endorsement of it. A projectible category licenses expectations that can be preserved, weakened, or defeated when conditions change. That gives kind claims a world-answerable dimension even when they remain plural and field-relative.

\section{The dispensable systemic-functional dependency}

The function-first conclusion doesn't require systemic-functional linguistics. A philosopher can reach the same result from easy ontology, modal normativism, inferentialism, pragmatism, or conceptual engineering. The needed premise is functional pluralism about discourse: not every term or sentence has the job of naming or describing an entity, property, or deep metaphysical structure.

Whatever its independent merits, SFL isn't needed for the metaphysical conclusion. Routing the argument through it can make a general point about semantic and inferential roles look dependent on a controversial grammatical architecture. Thomasson's own examples don't require that dependency. Application conditions, world-deferential rules, non-descriptive modal vocabulary, and pragmatic conceptual engineering can all be stated without taking on the additional commitments of systemic-functional grammar.

A small linguistic example makes the level point vivid. A semanticist may treat \mention{London} as a \term{proper name}: an expression used to refer to an individual place and to support inferences about reference, opacity, and identity. A syntactician may treat the same string as a \term{proper noun}: a nominal item whose distribution concerns noun-phrase position, determiner resistance, agreement, and modification. The referential role and the syntactic role are related, but neither can be read off from the other.

For this paper, the point is limited. I am not arguing that the linguistic framing is useless. It may help Thomasson organize the view and connect it to a broader account of discourse. The claim is narrower: the metaphysical payoff comes from function-first pluralism, not from function-in-syntax.

\section{The unmet opponent}

Thomasson's main natural-kind target is the essentialist tradition associated with Kripke and Putnam \citep{Kripke1980,Putnam1975Meaning}. The target is important, but it isn't the whole debate. The anti-essentialist projectibility tradition isn't treated as the central opponent.

The omission matters because the anti-essentialist realist tradition grants much of what Thomasson wants. Goodman made projectibility central to induction: not every true predicate supports reliable projection \citep{Goodman1955}. His point can be put without the technical details of \mention{grue}: two predicates can fit all observed cases while only one supports new expectations. The difference lies in projectibility, not in past truth alone.

Dupré rejects the idea that objects generally have one context-independent natural kind and defends a pluralist realism on which overlapping categories can be equally real \citep[5--7, 57--58, 262]{dupre1993disorder}. Boyd's homeostatic-property-cluster account, Khalidi's causal-network account, Magnus's pluralist critique of HPC reduction, and Slater's \term{natural kindness} all weaken or reject essentialism while retaining a realist role for kinds \citep{boyd1991,boyd1999,khalidi2013,khalidi2018nodes,magnus2012scientific,slater2015}.

Boyd gives the strongest bridge to Thomasson. In accommodationist terms, the reality of a kind turns on the contribution reference to it makes, within a disciplinary matrix, to reliable inference and accommodation to causal structure. Different species concepts can answer to different explanatory projects without becoming merely unreal or arbitrary \citep[169--176]{boyd1999}. Khalidi makes the parallel point in causal-network terms. Categories needn't be mind-independent; they need to serve epistemic purposes by tracking causal patterns and supporting projectible inferences \citep[202--203]{khalidi2013}.

Dupré still matters because he states the pluralist lesson bluntly. Category choice depends on the goals of inquiry, but the resulting kinds needn't be unreal. Species may be treated as individuals in some contexts and as modest kinds in others \citep[235]{dupre1993disorder}. He uses Goodmanian projectibility only in passing \citep[279 n.~20]{dupre1993disorder}, so he doesn't provide the full apparatus of support grades, demotion conditions, and field discipline that a projectibility-first account needs.

\section{Projectibility-first, field-relative kinds}

A projectibility-first account begins with the projection to be licensed. A category earns standing when membership in it supports reliable expectations about other features, behaviours, failure modes, or intervention outcomes. The first question isn't what essence all members share. The issue is what inferences the category licenses, under what conditions, and with what known failure modes.

This requires three distinctions. First, a category is a sorting device. Second, a profile is the pattern of properties, relations, or behaviours that makes the sorting projectible. Third, a kind claim says that the category earns standing by tracking that profile well enough for some field of inquiry or practice. Confusing these three levels makes weak categories look like well-supported kinds.

Field-relativity here isn't whim-relativity. A field is a stabilized practice or inquiry regime with its own questions, measurement routines, uptake patterns, authorities, records, and defeaters. Field-relative kind claims can be answerable to the world because their projections can fail. Categories can also be demoted: a category may remain useful for one task while losing standing for another.

The apparatus has three parts. Support grades concern the reliability, scope, and counterfactual robustness of the projections a category licenses. Failure conditions specify what would count as breakdown rather than noise: anomalous cases, field changes, intervention effects, or loss of explanatory integration. Demotion rules specify when a category remains useful for sorting or communication while losing kind-standing for a particular inquiry.

The \mention{fish} case shows the distinction. In ordinary food and market practices, \mention{fish} can be a useful category because it supports expectations about handling, preparation, and labelling. In phylogenetic systematics, the same category can be demoted because it doesn't track the ancestry relations that field asks about. In ecological work, a narrower category built around aquatic life may still support expectations about habitat, trophic role, or exposure. The word alone doesn't settle the kind claim; the field and projection do.

The \mention{platypus} example shows the other direction. Thomasson's world-deferential application conditions leave room for experts to correct ordinary use. A person can know how to use \mention{platypus} without knowing which traits carry the biological projections. Mammalogy can still treat the term as tied to a profile that supports expectations about related traits and explains why some apparent counterexamples are errors, immature cases, or abnormal cases. That's stronger than a loose grouping, but still not a field-neutral essence.

If HPC is read as a claim about stabilizing mechanisms alone, it can obscure this point. This realism is broader than homeostatic mechanism realism. A stable profile, a network position, a sustaining mechanism, and corrective control can come apart. Some projectible profiles are control-secured. Others draw support from history, isolation, institutional uptake, or background network structure. Homeostatic stabilization is one route to support, not the criterion for kindhood.

\section{The residual disagreement}

The disagreement with Thomasson should be stated narrowly. It isn't about whether purposes matter. They do. It isn't about whether empirical facts constrain concept and category choice. They do. Thomasson's own discussion of \mention{fish} shows that conceptual engineering can be worldly and non-arbitrary \citep[216--217]{thomasson2025rethink}.

Here the field-relative worry is that ordinary linguistic practice can appear to occupy the unmarked position. It's one field among others. Hard ontology is also a field, with its own vocabulary, questions, and internal standards. Its distance from ordinary use isn't by itself a defect. The defect would be failure to license projections, explanations, or disciplined contrasts that ordinary practice can't supply.

The dilemma gives hard ontology no free pass. If hard ontology claims to use ordinary terms, ordinary application conditions constrain it. If it introduces field-specific terms or category schemes, ordinary use no longer settles their standing, but the field owes projections, failure conditions, and demotion rules.

Nor is the worry that Thomasson preserves ordinary language against revision. Her conceptual engineering is explicitly revisionary. The issue is that, once ordinary linguistic practice, science, and metaphysics are all treated as fields, usefulness within one field doesn't by itself settle standing in another. A category can be useful, even indispensable, without yet earning kind-standing for a field-relative inquiry.

Projectibility may still add a realist constraint beyond empirically informed usefulness. A projectibility profile licenses defeasible expectations. It also yields tests: perturb a stabilizing relation, change a field condition, or move the category to a new setting, and predictive reliability should degrade in determinate ways. These aren't merely reports that a category serves our purposes. They're world-answerable claims: the world answers them by preserving or defeating the projected inferences.

Thomasson can respond that this still doesn't identify a field-neutral structure that mandates one correct category scheme. That response is right. The projectibility-first account doesn't need such a structure. It needs only the weaker realist claim that, once a field and projection are specified, the world can make some category profiles better supported than others.

\section{Conclusion}

Thomasson's function-first metaphysics should be accepted in its strongest form. It doesn't need systemic-functional linguistics to redirect metaphysics toward semantic rules, conceptual functions, and category choice. It also shouldn't be rejected by pretending that it can't handle natural-kind deference or empirical constraint. But ordinary use isn't a field-neutral court over every specialized practice, including metaphysics itself.

The anti-essentialist realist tradition remains largely outside her target. Projectibility-first realism accepts pluralism, rejects essences, and treats kind claims as field-relative. It then asks for support grades, failure conditions, and demotion rules. Thomasson's account hasn't yet displaced that extra discipline.

Projectibility-first realism doesn't ask for a privileged category scheme fixed independently of inquiry. It asks what follows once a field, a category, and a projection are specified. Some categories then earn more than usefulness: they support defeasible expectations, explain failures, and generate tests that can be preserved or defeated by the world. That's a realist constraint, but not an essentialist one.

\section*{Acknowledgements}

OpenAI Codex (GPT-5) served as a drafting and editing aid in the preparation of this paper. I am responsible for all theoretical claims, arguments, errors, and interpretive choices.

\clearpage
\printbibliography

\end{document}
